
\documentclass[a4paper,11pt]{article}
\usepackage[a4paper, margin=8em]{geometry}

% usa i pacchetti per la scrittura in italiano
\usepackage[french,italian]{babel}
\usepackage[T1]{fontenc}
\usepackage[utf8]{inputenc}
\frenchspacing 

% usa i pacchetti per la formattazione matematica
\usepackage{amsmath, amssymb, amsthm, amsfonts}

% usa altri pacchetti
\usepackage{gensymb}
\usepackage{hyperref}
\usepackage{standalone}

% imposta il titolo
\title{Appunti Comunicazioni Numeriche}
\author{Luca Seggiani}
\date{2026}

% imposta lo stile
% usa helvetica
\usepackage[scaled]{helvet}
% usa palatino
\usepackage{palatino}
% usa un font monospazio guardabile
\usepackage{lmodern}

\renewcommand{\rmdefault}{ppl}
\renewcommand{\sfdefault}{phv}
\renewcommand{\ttdefault}{lmtt}

% disponi il titolo
\makeatletter
\renewcommand{\maketitle} {
	\begin{center} 
		\begin{minipage}[t]{.8\textwidth}
			\textsf{\huge\bfseries \@title} 
		\end{minipage}%
		\begin{minipage}[t]{.2\textwidth}
			\raggedleft \vspace{-1.65em}
			\textsf{\small \@author} \vfill
			\textsf{\small \@date}
		\end{minipage}
		\par
	\end{center}

	\thispagestyle{empty}
	\pagestyle{fancy}
}
\makeatother

% disponi teoremi
\usepackage{tcolorbox}
\newtcolorbox[auto counter, number within=section]{theorem}[2][]{%
	colback=blue!10, 
	colframe=blue!40!black, 
	sharp corners=northwest,
	fonttitle=\sffamily\bfseries, 
	title=~\thetcbcounter: #2, 
	#1
}

% disponi definizioni
\newtcolorbox[auto counter, number within=section]{definition}[2][]{%
	colback=red!10,
	colframe=red!40!black,
	sharp corners=northwest,
	fonttitle=\sffamily\bfseries,
	title=~\thetcbcounter: #2,
	#1
}

% U.D.M
\newcommand{\amp}{\ensuremath{\mathrm{A}}}
\newcommand{\volt}{\ensuremath{\mathrm{V}}}
\newcommand{\meter}{\ensuremath{\mathrm{m}}}
\newcommand{\second}{\ensuremath{\mathrm{s}}}
\newcommand{\farad}{\ensuremath{\mathrm{F}}}
\newcommand{\henry}{\ensuremath{\mathrm{H}}}
\newcommand{\siemens}{\ensuremath{\mathrm{S}}}

% circuiti
\usepackage{circuitikz}
\usetikzlibrary{babel}

% disegni
\usepackage{pgfplots}
\pgfplotsset{width=10cm,compat=1.9}

% disponi codice
\usepackage{listings}
\usepackage[table]{xcolor}

\lstdefinestyle{codestyle}{
		backgroundcolor=\color{black!5}, 
		commentstyle=\color{codegreen},
		keywordstyle=\bfseries\color{magenta},
		numberstyle=\sffamily\tiny\color{black!60},
		stringstyle=\color{green!50!black},
		basicstyle=\ttfamily\footnotesize,
		breakatwhitespace=false,         
		breaklines=true,                 
		captionpos=b,                    
		keepspaces=true,                 
		numbers=left,                    
		numbersep=5pt,                  
		showspaces=false,                
		showstringspaces=false,
		showtabs=false,                  
		tabsize=2
}

\lstdefinestyle{shellstyle}{
		backgroundcolor=\color{black!5}, 
		basicstyle=\ttfamily\footnotesize\color{black}, 
		commentstyle=\color{black}, 
		keywordstyle=\color{black},
		numberstyle=\color{black!5},
		stringstyle=\color{black}, 
		showspaces=false,
		showstringspaces=false, 
		showtabs=false, 
		tabsize=2, 
		numbers=none, 
		breaklines=true
}

\lstdefinelanguage{javascript}{
	keywords={typeof, new, true, false, catch, function, return, null, catch, switch, var, if, in, while, do, else, case, break},
	keywordstyle=\color{blue}\bfseries,
	ndkeywords={class, export, boolean, throw, implements, import, this},
	ndkeywordstyle=\color{darkgray}\bfseries,
	identifierstyle=\color{black},
	sensitive=false,
	comment=[l]{//},
	morecomment=[s]{/*}{*/},
	commentstyle=\color{purple}\ttfamily,
	stringstyle=\color{red}\ttfamily,
	morestring=[b]',
	morestring=[b]"
}

% disponi sezioni
\usepackage{titlesec}

\titleformat{\section}
	{\sffamily\Large\bfseries} 
	{\thesection}{1em}{} 
\titleformat{\subsection}
	{\sffamily\large\bfseries}   
	{\thesubsection}{1em}{} 
\titleformat{\subsubsection}
	{\sffamily\normalsize\bfseries} 
	{\thesubsubsection}{1em}{}

% disponi alberi
\usepackage{forest}

\forestset{
	rectstyle/.style={
		for tree={rectangle,draw,font=\large\sffamily}
	},
	roundstyle/.style={
		for tree={circle,draw,font=\large}
	}
}

% disponi algoritmi
\usepackage{algorithm}
\usepackage{algorithmic}
\makeatletter
\renewcommand{\ALG@name}{Algoritmo}
\makeatother

% disponi numeri di pagina
\usepackage{fancyhdr}
\fancyhf{} 
\fancyfoot[L]{\sffamily{\thepage}}

\makeatletter
\fancyhead[L]{\raisebox{1ex}[0pt][0pt]{\sffamily{\@title \ \@date}}} 
\fancyhead[R]{\raisebox{1ex}[0pt][0pt]{\sffamily{\@author}}}
\makeatother

\begin{document}
% sezione (data)
\section{Lezione del 28-04-26}

% stili pagina
\thispagestyle{empty}
\pagestyle{fancy}

% testo
Avevamo visto in 18.4 la definizione di processo aleatorio, cioè sostanzialmente una "variabile aleatoria" che si traduceva in una \textit{funzione} di $t$ anziché un valore costante. In particolare, campionando un processo aleatorio ad un certo tempo $t^*$, ottenevamo effettivamente una variabile aleatoria come da definizione.

\subsubsection{Definizione statistica del primo ordine}
Abbiamo detto che fissato un istante $t$, $X(t)$ rappresenta una variabile aleatoria. La sua funzione di distribuzione è detta \textit{funzione di distribuzione del primo ordine} del processo $X(t)$:
$$
F_X(x; t) = \mathrm{Pr}\{X(t) \leq x\}
$$
Naturalmente questa funzione dipende anche da una variabile temporale perché le proprietà statistiche della variabile aleatoria $X(t)$ cambiano, in generale, al cambiare dell’istante di tempo $t$ al quale si "campiona" il processo. Chiaramente, si può anche descrivere una \textit{funzione densità di probabilità del primo ordine} del processo $X(t)$:
$$
f_X(x; t) = \frac{\partial F_X(x, t)}{\partial x}
$$

Per le distribuzioni del primo ordine si possono definire \textit{indici statistici} del primo ordine, cioè si può avere:
\begin{itemize}
\item La funzione \textit{valor medio} $\eta(x)$:
$$
\eta(x) = E\{X(t)\} = \int x f_X(x;t) \, dx
$$
\item La funzione \textit{potenza media statistica istantanea} $P_X(t)$:
$$
P_X(t) = E\{X^2(t)\} = \int x^2 f_X(x;t) \, dx
$$
\item La funzione \textit{varianza} $\sigma^2_X(t)$:
$$
\sigma^2_X(t) = E\left\{ (X(t) - \eta_X(t))^2 \right\} = \int (x - \eta_X(t)) f_X(x;t) \, dx = P_X(t) - \eta^2_X(t)
$$
\end{itemize}

\subsubsection{Definizione statistica del secondo ordine}
Dopo la definizione statistica del primo ordine, segue chiaramente quella del secondo ordine, che fissa 2 anziché 1 istanti temporali, ottenendo una \textit{funzione distribuzione di probabilità del secondo ordine} che è effettivamente una distribuzione di probabilità \textit{congiunta}:
$$
F_X(x_1, x_2; t_1, t_2) = \mathrm{Pr} \{X(t_1) \leq x_1, X(t_2) \leq x_2\}
$$
Come prima, possiamo definire anche la \textit{funzione densità di probabilità del secondo ordine}:
$$
f_X(x_1, x_2; t_1, t_2) = \frac{\partial^2 F_X(x_1, x_2; t_1, t_2)}{\partial x_1 \partial x_2}
$$

Il ragionamento può essere esteso a piacere, ovvero la descrizione statistica di un processo è completa solo quando si è in grado di caratterizzare il comportamento statistico congiunto di un numero $N$ arbitrario di variabili aleatorie $X(t_1), ..., X(t_N)$ estratte da $X(t)$ a $N$ istanti diversi, comunque grande sia l’intero $N$, e comunque si scelga la $N$-upla di istanti $\{t_1, ... t_N\}$. Questo processo è esattamente quello che avevamo definito \textit{descrizione statistica} di un processo aleatorio in 18.4.1.

\par\smallskip

\noindent
\textbf{\textsf{Autocorrelazione}}

\noindent
Al secondo ordine si possono definire degli \textit{indici statistici del secondo ordine} specifici, fra cui notiamo l'\textbf{autocorrelazione}. Fissiamo adesso due istanti di tempo arbitrari $t_1$ e $t_2$ sul nostro processo $X(t)$, ottenendo le due variabili aleatorie $Y = X(t_1)$ e $Z = X(t_2)$. A questo punto è significativa la \textit{correlazione}:
$$
r_{YZ} = E\left[YZ\right]
$$
fra queste due variabili.

Il valore di questa correlazione risulterà funzione dei due istanti $t_1$ e $t_2$ ai quali le variabili sono state estratte, e potrà essere calcolata solo conoscendo la funzione densità di probabilità del secondo ordine del processo:
$$
R_X(t_1, t_2) = E\left[ X(t_1) X(t_2) \right] = \int_{x_1 = -\infty}^{\infty} \int_{x_2 = -\infty}^{\infty} x_1 x_2 f_X(x_1, x_2; t_1, t_2) \, dx_1 dx_2
$$
La funzione $R_X(t_1, t_2)$ si chiama funzione di \textit{autocorrelazione} perché le due variabili aleatorie di cui si calcola la correlazione sono estratte dallo stesso processo aleatorio.

\par\smallskip

\noindent
\textbf{\textsf{Autocovarianza}}

\noindent
Se invece tra le due variabili aleatorie $Y$ e $Z$ calcoliamo la covarianza otteniamo la funzione di \textbf{autocovarianza} $C_X(t_1, t_2)$ di $X(t)$:
$$
C_X(t_1, t_2) = E\left[ (X(t_1) - \eta_X(t_1)) \, (X(t_2) - \eta_X(t_2)) \right]
$$
$$
= \int_{x_1 = -\infty}^{\infty} \int_{x_2 = -\infty}^{\infty} (X(t_1) - \eta_X(t_1)) \, (X(t_2) - \eta_X(t_2)) \, f_X(x_1, x_2; t_1, t_2) \, dx_1 dx_2
$$
$$
= R_X(t_1, t_2) - \eta_X(t_1) \eta_X(t_2)
$$

\par\smallskip

Notiamo poi che ponendo $t_1 = t_2 = t$, cioè valutando un unico punto, l'autocorrelazione e l'autocovarianza si identificano rispettivamente con il valore quadratico medio (potenza media statistica istantanea) e la varianza della variabile aleatoria $X(t)$:
$$
R_X(t, t) = E\left[ X^2(t) \right] = P_X(t), \quad C_X(t, t) = E\left[ (X(t) - \eta_X(t))^2 \right] = \sigma^2_X(t)
$$

\subsubsection{Descrizione in potenza}
In molti casi, ci si accontenta di studiare il processo analizzando solamente la funzione valore medio e la funzione di autocorrelazione di un processo aleatorio, cioè di effettuare la cosiddetta \textit{descrizione in potenza} del processo.

La funzione valore medio di un processo $X(t)$ è il valore medio della variabile $X(t)$. Solitamente questa è definita come funzione del tempo:
$$
\eta_X(t) = E\left[ X(t) \right] =
\begin{cases}
	\sum_i x_i \, \mathrm{Pr} \left\{ X(t) = x_i \right\} \\
	\int_{-\infty}^{\infty} x \, f_X(x; t) \, dx
\end{cases}
$$

La funzionedi autocorrelazione di un processo $X(t)$ è la correlazione (momento congiunto ordinario) delle variabili aleatorie $X(t_1)$ e $X(t_2)$. Essa è funzione di $t_1$ e $t_2$:
$$
R_X(t_1, t_2) = E \left[ X(t_1) X(t_2) \right] = 
\begin{cases}
	\sum_i \sum_j x_i x_j \, \mathrm{Pr} \left\{ X(t_1) = x_i , X(t_2) = x_j \right\}	\\
	\int_{x_1 = -\infty}^{\infty} \int_{x_2 = -\infty}^{\infty} x_1 x_2 f_X(x_1, x_2; t_1, t_2) \, dx_1 dx_2
\end{cases}
$$

\end{document}
